\subsection{线性形式；模的对偶}

设 \(\mathbf{E}\) 为一个左 \(\mathbf{A}\) -模。由于 \(\mathbf{A}\) 是一个 \((\mathbf{A}, \mathbf{A})\) -双模， \(\operatorname{Hom}_{\mathbf{A}}(\mathbf{E}, \mathbf{A}_s)\) 具有一个典范的右 \(\mathbf{A}\) -模结构（§1，第14号）。

定义2. 对于每个右A-模E，右A-模 \(\mathrm{Hom}_{\mathbf{A}}(\mathbf{E},\mathbf{A}_{s})\) 称为E的对偶模（或简称为E的对偶†），其元素称为E上的线性形式。

若 E 是一个右 A-模，则集合 \(\operatorname{Hom}_{\mathbf{A}}(\mathbf{E}, \mathbf{A}_{d})\) 连同其典范的左A-模结构同样称为E的对偶，其元素称为E上的线性形式。

在本章中，E* 将用于表示（左或右）A-模 E 的对偶。

例。*在 \(\mathbf{R}\) 的区间 \([a, b]\) 上的连续实值函数的向量空间（关于域 \(\mathbf{R}\) ）上，映射 \(x \mapsto \int_{a}^{b} x(t) dt\) 是一个线性形式。*

设 E 为左 A-模， \(E^{*}\) 为其对偶；对于每一对有序元素 \(x \in E\) , \(x^{*} \in E^{*}\) ，A 的元素 \(x^{*}(x)\) 记为 \(\langle x, x^{*} \rangle\) 。则关系式

\begin{enumerate}
\def\labelenumi{(\arabic{enumi})}
\setcounter{enumi}{8}
\tightlist
\item
\end{enumerate}

\[
\langle x + y, x ^ {*} \rangle = \langle x, x ^ {*} \rangle + \langle y, x ^ {*} \rangle\tag{10}
\]

\[
\langle x, x ^ {*} + y ^ {*} \rangle = \langle x, x ^ {*} \rangle + \langle x, y ^ {*} \rangle\tag{11}
\]

\[
\langle \alpha x, x ^ {*} \rangle = \alpha \langle x, x ^ {*} \rangle\tag{12}
\]

\[
\langle x, x ^ {*} \alpha \rangle = \langle x, x ^ {*} \rangle \alpha
\]

对 \(x, y\) 、 \(\mathbf{E}, x^*, y^*\) （其中x, y属于E，而x*, y*属于 \(\mathbf{E}^*\) ）和 \(\alpha \in \mathbf{A}\) 成立。映射 \((x, x^*) \mapsto \langle x, x^* \rangle\) （从 \(\mathbf{E} \times \mathbf{E}^*\) 到 \(\mathbf{A}\) 中）称为 \(\mathbf{E} \times \mathbf{E}^*\) 上的典范双线性型（双线性型的概念将在第九章，§1中一般地定义）。每个线性形式 \(x^*\) （在 \(\mathbf{E}\) 上）都可以视为对应于典范双线性型的部分映射 \(x \mapsto \langle x, x^* \rangle\) 。

当E是右A-模时， \(x^{*}(x)\) （线性形式 \(x^{*} \in E^{*}\) 在元素 \(x \in E\) 处的值）记为 \(\langle x^{*}, x \rangle\) ，与(11)和(12)对应的公式写为

\[
\begin{array}{l} \langle x ^ {*}, x \alpha \rangle = \langle x ^ {*}, x \rangle \alpha \\ \langle \alpha x ^ {*}, x \rangle = \alpha \langle x ^ {*}, x \rangle . \end{array}
\]

当A为交换环时，两种记号均可使用。

对于每个环A，将每个 \(\xi\in\mathbf{A}\) 对应到线性形式 \(\eta\mapsto\eta\xi\) （ \(\mathbf{A}_{s}\) 上的）的映射，是 \(\mathbf{A}_{d}\) 到 \(\mathbf{A}_{s}\) 的对偶上的一个同构。

这是典范同构 \(E \to \operatorname{Hom}_{A}(A_{s}, E)\) （§1，第14号，注记2）对应于 \(E = A_{s}\) 的特殊情形；§1 第14号中图表（50）的交换性表明，我们在这里得到的是右A-模之间的一个同构。

若利用命题5中的同构把 \(A_{d}\) 与 \(A_{s}\) 的对偶等同起来，则 \(A_{s} \times A_{d}\) 上的典范双线性形式由

\[
\langle \xi , \xi^ {*} \rangle = \xi \xi^ {*} \quad \text { for } \xi , \xi^ {*} \text { in } A.\tag{13}
\]

类似地， \(A_{d}\) 被典范地等同于 \(A_{s}\) 的对偶， \(A_{d} \times A_{s}\) 上的典范双线性形式由

\[
\langle \xi^ {*}, \xi \rangle = \xi^ {*} \xi \quad \text { for } \xi , \xi^ {*} \text { in } A.\tag{14}
\]

